\begin{afterword}
\summaryhead

The post asks why we expect the Sun to rise tomorrow. Past sunrises, the laws of physics and Occam's razor all support the expectation only if the future resembles the past. Some people conclude that science rests on faith, and use this to shield their own beliefs from criticism. The author calls that a bad sign, but grants the real problem: if every belief needs a justification, how does the chain end, and if it may end in an unjustified assumption, why not assume anything?

Bayesian updating does not solve this, since the answer depends on the prior, and possible minds with anti-Occamian priors defend them by their own rule. The author's answer is to examine questions like ``Should I trust my brain?'' with the current mind, since there is nothing else to use. The chain of justification then loops: induction has worked, simplicity explains why, and evolution explains why the brain can see it. The author holds such reflective loops to be different from circular logic, such as a paper that declares itself true. The Bible's claim to be God's word fails on its record, not because it loops. The aim is not to persuade an ideal philosopher of perfect emptiness but to win, questioning everything, including the rules of reasoning, with full force.

\responsehead

The post states the problem fairly. It grants that its earlier reply to ``Science is based on faith too'' did not answer the philosophical problem. It concedes that Bayesian updating does not solve induction. And it raises two strong objections to its own view before answering them: a counterinductive mind can defend its rule by that rule, and most minds that settle into coherence would be incorrect. Its answer, to examine one's methods with one's current mind and to trust the loop because of a causal story about how the mind was made, is a coherent position.

The problem and much of the answer are old. The sunrise is Hume's example. The regress goes back to Aristotle. The image of repairing beliefs without starting from scratch is Neurath's boat. The evolutionary argument is Quine's. The idea that rules of induction are justified by mutual adjustment with accepted inferences is Goodman's, and the reply that rationality can be held open to its own criticism is Bartley's. Two days later the author wrote that others had considered the question before. This is information for the reader, and most of it supports the post.

The weakest point is the one the post rests on: that reflective loops can be told from circular logic ``by common sense''. The post suggests tests along the way, a causal story of how the mind tracks the world and a record on other matters, but it does not state a criterion, and the author said two days later that the difference had not been formalized. The literature offered a candidate, rule-circular against premise-circular arguments, which fits the post's two examples.

One limit of scope. The post's dilemma, an endless chain or an unjustified stopping point, leaves out the classic answer that some beliefs are justified without resting on others; for induction that would mean an a priori justification, which the author rejected elsewhere.

In short: a fair statement of the regress problem and a coherent naturalist answer with a long pedigree (Hume, Neurath, Quine, Goodman), whose central distinction between reflective loops and circular logic is left to common sense.
\end{afterword}
